\documentclass[11pt]{article}
\usepackage[margin=1in]{geometry}
\usepackage{amsmath,amssymb,amsthm,mathtools}
\usepackage{enumitem}
\usepackage{hyperref}
\hypersetup{colorlinks=true,linkcolor=blue,citecolor=blue,urlcolor=blue}

\newtheorem{theorem}{Theorem}
\newtheorem{proposition}{Proposition}
\newtheorem{lemma}{Lemma}
\newtheorem{remark}{Remark}
\newtheorem{question}{Question}

\newcommand{\dd}{\,d}
\newcommand{\one}{\mathbf 1}
\newcommand{\T}{T_W}
\newcommand{\U}{U}
\newcommand{\DeltaTwo}{\Delta_2}

\title{The directed weighted transitivity route for the smoothed Goodman inequality}
\author{Working note}
\date{\today}

\begin{document}
\maketitle

\section{Purpose of this note}

This note records a focused attempt to prove the smoothed Goodman inequality by the
``directed weighted transitivity / degree-bias'' route.  The target inequality is
\begin{equation}
\DeltaTwo(W)
:=t(\theta_{1,2,4},W)-(2p-1)t(C_5,W)\ge 0,
\qquad p=t(K_2,W),
\tag{SG2}
\end{equation}
where \(\theta_{1,2,4}=C_3\cup_{K_2}C_5\) is the theta graph with three internally
vertex-disjoint paths of lengths \(1,2,4\) between the same two endpoints.
Equivalently,
\begin{equation}
\DeltaTwo(W)
=\iint W(x,y)T_W^2(x,y)T_W^4(x,y)\dd x\dd y
 -(2p-1)\iint W(x,y)T_W^4(x,y)\dd x\dd y .
\end{equation}
For \(p\le 1/2\) the inequality is trivial.  The real case is \(p>1/2\).

The main outcome of this note is not a complete proof.  It is a cleaner global
reformulation:
\begin{quote}
\emph{The inequality \((\mathrm{SG2})\) is exactly an averaged weighted
transitivity inequality for the complement-neighbourhoods of vertices, where the
edge weight is the \(C_5\)-edge kernel \(B_W(x,y)=W(x,y)T_W^4(x,y)\).}
\end{quote}
This identifies the precise compensation mechanism.  A tempting pointwise
version is false, so any proof must exploit averaging over all complement rows
or use a flag/SOS certificate.

\section{The degree-bias decomposition}

Let
\[
U:=1-W,\qquad q:=1-p,
\qquad d(x):=\int W(x,y)\dd y,
\qquad u(x):=\int U(x,y)\dd y=1-d(x).
\]
Define
\begin{align}
G(W)&:=\iint W(x,y)T_U^2(x,y)T_W^4(x,y)\dd x\dd y,\label{eq:G}\\
D(W)&:=\int (p-d(x))T_W^5(x,x)\dd x
      =\int (u(x)-q)T_W^5(x,x)\dd x.\label{eq:D}
\end{align}
The elementary identity
\begin{equation}
T_W^2(x,y)-(2p-1)
=T_U^2(x,y)-(p-d(x))-(p-d(y))
\label{eq:basic-identity}
\end{equation}
gives
\begin{equation}
\boxed{\DeltaTwo(W)=G(W)-2D(W).}
\label{eq:Gminus2D}
\end{equation}
Thus the global theorem is equivalent to
\begin{equation}
\boxed{G(W)\ge 2D(W).}
\tag{TD}
\end{equation}
The interpretation is as follows.  The term \(D(W)\) is positive precisely when
rooted \(C_5\)-mass is biased toward vertices of below-average degree.  The term
\(G(W)\) counts, with the same \(C_5\)-edge weight, edges whose two endpoints have
a common non-neighbour.  Thus \((\mathrm{TD})\) says that every possible
low-degree bias of rooted \(C_5\)'s is compensated by a weighted transitivity
defect in the complement relation.

\section{The \(C_5\)-edge kernel and the averaged transitivity formula}

Define the \(C_5\)-edge kernel
\begin{equation}
B_W(x,y):=W(x,y)T_W^4(x,y).
\end{equation}
Let
\begin{equation}
Z_W:=\iint B_W(x,y)\dd x\dd y=t(C_5,W),
\qquad
R_W(x):=\int B_W(x,y)\dd y=T_W^5(x,x).
\end{equation}
For each vertex \(z\), let
\begin{equation}
f_z(x):=U(z,x),
\qquad
u(z):=\int f_z(x)\dd x=u(z).
\end{equation}
Think of \(f_z\) as the complement-neighbourhood profile of \(z\).
Now define the one-vertex weighted transitivity defect
\begin{equation}
\mathcal F_W(z):=
\iint B_W(x,y)f_z(x)f_z(y)\dd x\dd y
-2\int R_W(x)f_z(x)\dd x
+2\nu(z)Z_W.
\label{eq:Fz}
\end{equation}

\begin{proposition}[Averaged directed transitivity identity]
For every graphon \(W\),
\begin{equation}
\boxed{\DeltaTwo(W)=\int_0^1 \mathcal F_W(z)\dd z.}
\label{eq:avgF}
\end{equation}
Consequently, \((\mathrm{SG2})\) is equivalent to
\begin{equation}
\boxed{\int_0^1 \mathcal F_W(z)\dd z\ge0.}
\end{equation}
\end{proposition}

\begin{proof}
The first term in \(\int\mathcal F_W(z)\dd z\) is
\[
\iiint B_W(x,y)U(z,x)U(z,y)\dd x\dd y\dd z
=\iint B_W(x,y)T_U^2(x,y)\dd x\dd y=G(W).
\]
The second term is
\[
-2\iint R_W(x)U(z,x)\dd x\dd z
=-2\int R_W(x)u(x)\dd x.
\]
The third term is
\[
2Z_W\int u(z)\dd z=2qZ_W.
\]
Therefore
\[
\int \mathcal F_W(z)\dd z
=G(W)-2\int (u(x)-q)R_W(x)\dd x
=G(W)-2D(W)=\DeltaTwo(W).
\]
\end{proof}

\section{The coverage interpretation in the random-free case}

Assume for this paragraph that \(W\) is random-free, so \(W\in\{0,1\}\).  For a
fixed \(z\), let
\[
S_z:=\{x:U(z,x)=1\}
\]
be the non-neighbourhood of \(z\).  Sample an edge \((X,Y)\) from the \(C_5\)-edge kernel \(B_W\), i.e. with density proportional to
\(B_W(x,y)\).  Then
\[
\frac{1}{Z_W}\iint B_W(x,y)f_z(x)f_z(y)\dd x\dd y
=\mathbb P(X,Y\in S_z),
\]
and
\[
\frac{1}{Z_W}\int R_W(x)f_z(x)\dd x
=\mathbb P(X\in S_z).
\]
Hence
\begin{equation}
\frac{\mathcal F_W(z)}{Z_W}
=2\mu(S_z)-\mathbb P(\{X,Y\}\cap S_z\ne\varnothing),
\label{eq:coverage}
\end{equation}
where \(\mu(S_z)=u(z)\) is the ordinary Lebesgue mass of the non-neighbourhood.
Thus the global inequality says
\begin{equation}
\boxed{
\mathbb E_{z}\,\mathbb P_{(X,Y)\sim B_W}\bigl(\{X,Y\}\cap S_z\ne\varnothing\bigr)
\le 2q.}
\label{eq:coverage-global}
\end{equation}
This is the directed weighted transitivity form: under the \(C_5\)-edge-biased
measure, complement-neighbourhoods should not cover more edge mass than twice
the global complement density, once averaged over \(z\).

\section{A tempting pointwise strengthening is false}

The most tempting route would be to prove
\begin{equation}
\mathcal F_W(z)\ge0\qquad\text{for a.e. }z.
\tag{false-pointwise}
\end{equation}
Equivalently, in the random-free case, every individual non-neighbourhood
\(S_z\) would cover at most \(2\mu(S_z)\) of the \(C_5\)-edge-biased edge mass.
This is false, even on the top branch.

Consider the equal four-block random-free graphon with kernel-value matrix
\begin{equation}
A=
\begin{pmatrix}
0&1&1&1\\
1&1&1&0\\
1&1&0&0\\
1&0&0&0
\end{pmatrix}.
\end{equation}
Here
\[
p=\frac{9}{16},\qquad Z_W=t(C_5,W)=\frac{3}{32}.
\]
A direct exact calculation gives
\[
\DeltaTwo(W)=\frac{33}{1024}>0,
\]
but the four values of \(\mathcal F_W(z)\) are
\[
-\frac1{256},\qquad \frac{21}{512},\qquad \frac{23}{512},
\qquad \frac3{64}.
\]
Thus one vertex has negative pointwise transitivity defect, while the averaged
defect is positive.  The proof must therefore use a genuinely global averaging
mechanism.  The accompanying script \texttt{verify\_directed\_transitivity.py}
checks these identities exactly.

\section{What this route has achieved, and what remains}

The directed transitivity route has not yet produced the full theorem, but it
has replaced the previous somewhat opaque inequality \(G\ge2D\) by the following
precise problem.

\begin{question}[Averaged weighted transitivity]
For every graphon \(W\), with \(B_W=W T_W^4\), is it true that
\begin{equation}
\int_0^1
\left[
\iint B_W(x,y)f_z(x)f_z(y)\dd x\dd y
-2\int R_W(x)f_z(x)\dd x
+2u(z)Z_W
\right]
\dd z\ge0?
\end{equation}
Equivalently, in the random-free case, does the average \(C_5\)-edge coverage
of non-neighbourhoods satisfy \eqref{eq:coverage-global}?
\end{question}

This is exactly equivalent to \((\mathrm{SG2})\), but it is more informative than
the raw quantum graph expression.  It says that the only possible obstruction is
not merely low degree bias, but low degree bias that is not accompanied by enough
\(C_5\)-weighted complement transitivity defect.

The pointwise counterexample above also clarifies the next proof choices:
\begin{enumerate}[label=(\roman*)]
\item A purely pointwise set-isoperimetric proof cannot work.
\item A successful human proof must exploit cancellation between different
vertices \(z\), or a structural property of the family of complement rows
\(f_z=U(z,\cdot)\).
\item A flag/SOS certificate is very natural here: the expression
\(\int\mathcal F_W(z)\dd z\) is a fixed seven-vertex quantum graph inequality,
while \(\mathcal F_W(z)\) itself is a one-vertex flagged expression over a
\(C_5\)-edge-biased type.
\end{enumerate}

\section{Recommended next move}

The directed transitivity formulation suggests that the next high-value attack
should be one of the following.

\begin{enumerate}[label=\textbf{Route \arabic*.},leftmargin=*]
\item \textbf{Averaging/cancellation proof.}  Try to prove directly that negative
values of \(\mathcal F_W(z)\) are compensated by positive values of
\(\mathcal F_W(z')\) along complement edges or along common complement
neighbourhoods.  The four-block example above is the smallest test case for any
such cancellation lemma.

\item \textbf{Flag/SOS proof in the flagged language.}  Encode
\(\int\mathcal F_W(z)\dd z\ge0\) as a flag-algebra inequality with one labelled
outside vertex and a \(C_5\)-edge-biased structure.  Since the pointwise flag is
not nonnegative, the certificate must combine several flags or use an
8-vertex lift/disconnected term to express the averaging over \(z\).

\item \textbf{Weighted transitivity theorem for complement rows.}  Search for a
statement of the form
\[
\int \mathcal F_W(z)\dd z
\ge c\,\mathcal T(W)
\]
where \(\mathcal T(W)\) is a nonnegative measure of failure of the complement
relation to be an equivalence relation.  Equality should occur exactly at
balanced complete multipartite graphons and at \(W\equiv1\).
\end{enumerate}

At the current stage, this route has clarified the mechanism but has not yet
closed the theorem.  The key lesson is that the desired inequality is an
\emph{averaged} directed transitivity statement, not a pointwise one.

\end{document}
